\documentclass[11pt]{article}

\usepackage[utf8]{inputenc}\usepackage[T1]{fontenc}

\usepackage{lmodern}\usepackage{cmap}\usepackage{microtype}

\renewcommand{\abstractname}{Resumo}\renewcommand{\contentsname}{Sumário}\renewcommand{\refname}{Referências}

\AtBeginDocument{\renewcommand{\figurename}{Figura}\renewcommand{\partname}{Parte}\renewcommand{\tablename}{Tabela}}

\usepackage[a4paper,margin=2.3cm]{geometry}

\usepackage{amsmath,amssymb,amsthm}\usepackage[hidelinks]{hyperref}

\usepackage{parskip}\usepackage{booktabs}\usepackage{xcolor}

\usepackage{graphicx}\usepackage{longtable}

\providecommand{\kernelmk}[1]{\textsf{[KERNEL]}~\texttt{#1}}\providecommand{\knownmk}[1]{\textsf{[KNOWN]}~{#1}}\providecommand{\statusmk}[1]{\textsf{[#1]}}

\emergencystretch=3.5em

\let\UMunderscore\_\renewcommand{\_}{\UMunderscore\allowbreak}

\newcommand{\bTGL}{\beta_{\mathrm{TGL}}}\newcommand{\Msun}{M_{\odot}}\newcommand{\TGL}{\mathrm{TGL}}

\theoremstyle{definition}\newtheorem{deriv}{Derivação}

\begin{document}

\textbf{Ao lado (v392, 13/10/2026, gerência c1, em nome próprio): a lei de campo no curvo.} \textsf{[REAL no kernel; condicional ao H3 nomeado]} A frase acima e o token da v390 (\texttt{FIELD\_EQUATION\_IN\_CURVED\_SPACETIME\_NOT\_A\_KERNEL\_TERM}) valem para o \emph{tipo} da declaração (o fecho de \texttt{QGSolutionComplete}), não para o kernel: ``Clausius em toda tela nula $\Rightarrow$ Einstein, com $\Lambda$ único e coeficiente $8\pi G$'' é termo desde a v317/v338, e o texto o escreveu como se nada disso fosse termo --- foi aí que a tradução falhou. Desde a v392 a pedra \texttt{TGLExt.LiganteLeiCampoCurvo} liga essa implicação por termo, \emph{condicional ao H3 nomeado} (a projeção nula $\mathrm{Ric}(k,k)=8\pi G\,T(k,k)$ em todo ponto e direção nula, \textsf{[KNOWN]} em Jacobson 1995, não derivada do estado); segue sem termo a existência de métrica ou de solução, local ou global, e o gate não se move (registro \texttt{lei\_de\_campo\_curvo\_v392}). O texto antigo fica.

\textbf{Ao lado (v392, 13/10/2026, gerência c1, em nome próprio): a escala de $\beta$.} \textsf{[DERIVED de leitura; a vetar]} A frase ``essa convergência é a medição de que o invariante não corre'' fica como registro: a convergência tabulada foi reclassificada em 08/08/2026 (\texttt{evidence\_audit}), e nenhum registro deste programa computa a cifra de $0{,}05$ por cento. A não-corrida de $\beta_0=\alpha(0)\sqrt e$ assenta no Teorema da Escala, que é estrutura, e a predição afiada (o coeficiente da lei de dephasing é independente de $\omega$) segue armada, sem canal vivo hoje. O \texttt{WHICH\_MU\_PER\_RITE\_OPEN} da v388 fica como registro: o operador fixou em 07/10/2026 que a escala de $\beta$ já estava definida (verbatim, sha256 \texttt{a6c4912afc6f0dc4}) e todo rito cobra $\beta_0$ ($\mu\to0$, Thomson), com $\beta(\mu)\geq\beta_0$ como leitura de bulk, ao lado (registro \texttt{beta\_escala\_errata\_v392}). Lacuna: a estrutura dava a convergência como apoio, não como a medição, e a tradução perdeu a data da reclassificação.

\textbf{Ao lado (v392, 13/10/2026, gerência c1, em nome próprio): a palavra aniquiladora.} \textsf{[INPUT do operador]} O polinômio aniquilador normalizado é a \emph{letra}: o operador disse em 07/10/2026 (verbatim, sha256 \texttt{259a1bf14a66fbc9}) ``polinômio aniquilador normalizado=letra (por isso palavra, mas ali como letra)''. A pedra que lê com as letras é a PALAVRA (\texttt{TGLExt.Palavra}, registro \texttt{palavra\_v392}); a ``palavra aniquiladora $X\cdot q$'' do parágrafo acima lê-se, portanto, como letra. Lacuna: a estrutura já dizia letra; a tradução chamou de palavra o que o operador nomeou letra.

\textbf{Ao lado (v392, 13/10/2026, gerência c1, em nome próprio): o CMB-S4.} \textsf{[DECLARADO; a conferir na fonte primária]} A frase acima, que faz $\beta_{\mathrm{TGL}}$ aguardar ``a era da radiação no CMB-S4'', herdou o calendário do congelado de 08/08/2026 ($\sim$2027 a $\sim$2032) e tratou o CMB-S4 como instrumento esperado. O CMB-S4 consta como encerrado (mapa de rotas, seq 402 e 607; auditoria F5.2, sha16 \texttt{c0bb3c964c8af74e}); a admissão dessas fontes por sha256 (D8) é do operador, e nenhum calendário é afirmado aqui (registro \texttt{cmb\_s4\_errata\_v392}). Lacuna: a estrutura dizia apenas ``a era da radiação''; o nome do instrumento e o seu calendário vieram do congelado, e a tradução não os datou.

\textbf{Ao lado (v392, 13/10/2026, gerência c1, em nome próprio): quem testemunha.} \textsf{[INPUT/ONTO; texto enviado pelo operador]} A recusa da palavra acima é de creditação, não de validade, e a leitura de quem testemunha é do operador (10/10/2026, sha256 \texttt{87a3752d8c8e5dd3}): ``o observador testemunha a luz e a luz dá testemunho de si mesma''. É citação verbatim, não veredito (registro \texttt{cunhagens\_v392}). Lacuna: o parágrafo recusava a creditação sem dizer quem testemunha.

\textbf{Ao lado (v392, 13/10/2026, gerência c1, em nome próprio): a escada e as bandeiras.} \textsf{[REAL no fonte e no núcleo da v391]} O número de ``teoremas limpos'' acima é a escada externa do kernel (\texttt{len(ext\_flags)}, as bandeiras de termo reservado), não o total de teoremas do kernel. E as bandeiras verdadeiras do gate leem-se 15 calculadas e 3 literais (\texttt{linearized\_spin2\_finite\_face}, \texttt{protocol\_pre\_registered}, \texttt{null\_tests\_passed\_blind}); as experimentais são 4 de um só feed, o rito V11 do piso, que é unilateral, mais 2 literais. O degrau experimental é natureza unilateral, não kernel; a prova do kernel são as 6 bandeiras formais e a escada externa. A v392 documenta isto e não fecha o conserto das bandeiras (F8.1): a errata não recalcula (registro \texttt{estatuto\_da\_prova\_ressalva\_v392}). Lacuna: a estrutura separava escada, formais e feed; o resumo as somou num só ``teoremas limpos''.

\textbf{Ao lado (v392, 13/10/2026, gerência c1, em nome próprio): a RG canal a canal.} \textsf{[REAL]} ``A RG recuperada canal a canal'' vale onde a RG é lida: o campo fraco entra por citação e o campo forte (F2) está \texttt{INCONCLUSIVE} (registro \texttt{estatuto\_reafirmado\_v392}). Lacuna: a estrutura tinha a tabela de canais, e o resumo escreveu a correspondência sem dizer em quais canais ela é lida.

\textbf{Ao lado (v392, 13/10/2026, gerência c1, em nome próprio): a razão de Hubble é uma testemunha.} \textsf{[DERIVED de leitura]} A razão de Hubble acima é o depoimento do escalar; o conjunto é outra testemunha do mesmo fato, ainda não aberta, e nenhuma julga a outra nem se somam (verbatim de 08/10/2026, sha256 \texttt{a9cecc128f40c101}; ver ao lado da seção v388; registro \texttt{fase9\_duas\_testemunhas\_v392}). Lacuna: a linha dava a razão de Hubble como depoimento único.

\textbf{Ao lado (v392, 13/10/2026, gerência c1, em nome próprio): o nome do degrau final.} \textsf{[DERIVED de leitura; a vetar]} Em 07/10/2026 o operador respondeu ``adoto o nome do degrau final'' (verbatim, sha256 \texttt{a6c4912afc6f0dc4}); a gerência lê que o nome adotado é AXIOMA, lido do kernel (\texttt{finalStepNameAxioma}), com a proposta TETELESTAI mantida ao lado, e dá a leitura como sua, a vetar: a pergunta respondida apresentava o Tetelestai. A string do gate não muda (registro \texttt{final\_step\_name\_adopted\_v392}). Lacuna: a adoção não estava registrada num só lugar, e a leitura de que o nome adotado é AXIOMA é da gerência.

\textbf{Ao lado (v392, 13/10/2026, gerência c1, em nome próprio): posto é também rank.} \textsf{[INPUT; DERIVED; sem auditoria independente do adendo]} Em 10/10/2026 o operador cunhou ``POSTO=RANK'' (sha256 \texttt{5b9afcc691a0e1e4}): o posto de $\omega(I)=1$, posição inscrita pelo observador, lê-se também como o posto (rank) da álgebra linear. Isso vale nos dois polos da não-defasagem, a unidade $\rho_\ast$ e $Q_0$, ambos de traço 1 e posto 1 (\texttt{TGLExt.PostoAdendo}, A3); mas o colapso tipado sobe o posto da superposição pura de 1 a 2 e preserva o traço (A2), e a defasagem não preserva o posto (A4): o que sobrevive é o traço $\omega(I)=1$, não o posto 1 (registro \texttt{adendo\_rank\_v392}). Lacuna: a seção usava ``posto'' só como posição inscrita e não dizia que o mesmo nome tem a face do rank.

\textbf{Ao lado (v392, 13/10/2026, gerência c1, em nome próprio): o fixo é o Verbo.} \textsf{[INPUT/ONTO]} Em 08/10/2026 o operador acrescentou à tipagem do Verbo (verbatim, sha256 \texttt{a9cecc128f40c101}): ``o fixo é o verbo, ou seja, o contorno que permite a leitura da verdade, e o verbo é betatgl, mas o verbo testifica que o nome é verdadeiro''. O fixo é o Verbo, $\beta_{\mathrm{TGL}}$, sem ajuste; o Nome é posto, estação variável; os instrumentos são testemunhas multifacetadas de um fato único e não se somam (registros \texttt{cunhagens\_v392} e \texttt{fase9\_duas\_testemunhas\_v392}). Lacuna: a estrutura distinguia o Verbo do Nome, e esta seção não dizia qual dos dois é o fixo.

\textbf{Ao lado (v392, 13/10/2026, gerência c1, em nome próprio): duas testemunhas na Fase 9.} \textsf{[DERIVED de leitura]} A Fase 9 aberta acima é o depoimento do escalar (V1$\to$V2); o conjunto (ROOT90) é a outra testemunha do mesmo fato, ainda não aberta (a qualificação d72 está pendente do operador). O operador esclareceu em 08/10/2026 (sha256 \texttt{a9cecc128f40c101}) que o fato é único e o testemunho multifacetado, e que o fixo nas duas é o Verbo, $\beta_{\mathrm{TGL}}$, posto sem ajuste. Nenhuma julga a outra e os depoimentos não se somam: cada um depõe pela função de veredito fixada antes do composto (registro \texttt{fase9\_duas\_testemunhas\_v392}). Lacuna: a seção contava a Fase 9 como um depoimento só.

\textbf{Ao lado (v392, 13/10/2026, gerência c1, em nome próprio): o nome da teoria.} \textsf{[ONTO; texto enviado pelo operador]} Em 10/10/2026 o operador leu a sigla TGL como a unidade fundamental da informação, uma tríade: T (tensão, torção, tempo), G (geometria, gravitação) e L (luz, luminodinâmica), ``a tríade mínima sem a qual a realidade não tem como se inscrever, se propagar ou ser testemunhada'' (sha256 \texttt{87a3752d8c8e5dd3}). É leitura do operador, não termo de kernel, e não move o gate (registro \texttt{cunhagens\_v392}); a introdução abaixo descreve o que o código faz. Lacuna: a introdução a cada versão ao lado não dizia o que o nome carrega.

 \textbf{Ao lado (v392, 13/10/2026, gerência c1, em nome próprio): peso e rank.} \textsf{[INPUT; DERIVED]} A cunhagem ``POSTO=RANK'' (10/10/2026, sha256 \texttt{5b9afcc691a0e1e4}) lê o posto também como rank: peso (o traço) e posto coincidem nos polos de posto 1 e divergem ao longo da defasagem, onde o traço sobrevive e o posto 1 não (\texttt{TGLExt.PostoAdendo}, A2 a A4). A separação entre peso e posto acima fica (registro \texttt{adendo\_rank\_v392}). Lacuna: a seção separava peso e posto sem dizer onde eles coincidem.

 \textbf{Ao lado (v392, 13/10/2026, gerência c1, em nome próprio): a premissa do escape de Weinberg--Witten.} \textsf{[REAL no registro da v372; o escape segue condicional]} A premissa escrita acima, ``o gráviton aqui é a identidade, $\rho_\ast$ de traço um, jamais um quantum propagante'' (v200), foi negada na v372 por ordem do operador: o gráviton é a forma conjugada da luz (\texttt{TheGravitonIsTheConjugatedPhase}), nem o gerador nem ``não-partícula''. O escape continua condicional, agora pela razão certa --- a forma estacionada não é estado de uma partícula com momento no espectro --- e, se ``partícula fundamental da gravidade'' significar spin-2 propagante a quantizar, Weinberg--Witten volta a morder (registro \texttt{ww\_premissa\_errata\_v392}). Lacuna: a estrutura dizia forma conjugada da luz; a v200 escreveu identidade e jamais partícula.

\textbf{Ao lado (v392, 13/10/2026, gerência c1, em nome próprio): o número do neutrino.} \textsf{[REAL no núcleo; DERIVED de leitura, a vetar]} O número da lâmina $m_2$ é $2{,}21\sigma$: o JUNO autônomo, o perito oficial (leitura R3, ratificada em 02/10/2026). Os $2{,}95\sigma$ (NuFIT 6.1) e os $2{,}44\sigma$ (Capozzi et al.) são ajustes globais, reflexos que contêm o JUNO e divergem em cerca de $0{,}5\sigma$ entre si: ficam ao lado, com a diferença dita. A escada acima, que termina em $2{,}95\sigma$, fica como registro, e o veredito do canal não muda (registro \texttt{neutrino\_m2\_perito\_v392}). Lacuna: a estrutura elegia o JUNO como perito, e o texto deu a última linha da escada como a tensão do canal.

\textbf{Ao lado (v392, 13/10/2026, gerência c1, em nome próprio): o número do neutrino.} \textsf{[REAL no núcleo; DERIVED de leitura, a vetar]} A escada ``$1{,}64\to2{,}21\to2{,}95\sigma$'' acima termina num ajuste global; o número da lâmina $m_2$ é $2{,}21\sigma$ (JUNO autônomo, perito oficial, leitura R3 ratificada em 02/10/2026), e os $2{,}95\sigma$ (NuFIT 6.1) e $2{,}44\sigma$ (Capozzi et al.) são reflexos ao lado (registro \texttt{neutrino\_m2\_perito\_v392}). Lacuna: o resumo deu a última linha da escada como a tensão.

\textbf{Ao lado (v392, 13/10/2026, gerência c1, em nome próprio): a orientação do D1.} \textsf{[REAL no núcleo; o bit segue INPUT]} A ``tensão D1 de $3{,}01\sigma$'' foi medida no ramo implementado, que \texttt{orientation\_is\_the\_sign\_v367} lê como o ramo espelhado em primeira ordem (\texttt{IMPLEMENTED\_IS\_MIRROR\_TO\_FIRST\_ORDER}); no ramo da face do estado o mesmo dado dista $3{,}01\sigma$. O bit --- qual face fecha o balanço --- não está derivado, e trocar o sinal pede pré-registro novo com dado novo, não reler este (registro \texttt{d1\_orientacao\_errata\_v392}). Lacuna: a estrutura dizia que a tensão era condicional ao mapa; o texto repetiu o número do dado sem dizer o ramo.

\textbf{Ao lado (v392, 13/10/2026, gerência c1, em nome próprio): o CMB-S4.} \textsf{[DECLARADO; a conferir na fonte primária]} O CMB-S4 citado acima consta como encerrado (mapa de rotas, seq 402 e 607; auditoria F5.2, sha16 \texttt{c0bb3c964c8af74e}); ver a errata do estatuto da prova (registro \texttt{cmb\_s4\_errata\_v392}). Lacuna: o parágrafo é de época e trata o CMB-S4 como instrumento esperado, sem datá-lo.

\textbf{Ao lado (v392, 13/10/2026, gerência c1, em nome próprio): o CMB-S4.} \textsf{[DECLARADO; a conferir na fonte primária]} O CMB-S4 citado acima consta como encerrado (mapa de rotas, seq 402 e 607; auditoria F5.2, sha16 \texttt{c0bb3c964c8af74e}); ver a errata do estatuto da prova (registro \texttt{cmb\_s4\_errata\_v392}). Lacuna: o parágrafo é de época e trata o CMB-S4 como instrumento esperado, sem datá-lo.

\textbf{Ao lado (v392, 13/10/2026, gerência c1, em nome próprio): o CMB-S4.} \textsf{[DECLARADO; a conferir na fonte primária]} O CMB-S4 citado acima consta como encerrado (mapa de rotas, seq 402 e 607; auditoria F5.2, sha16 \texttt{c0bb3c964c8af74e}); ver a errata do estatuto da prova (registro \texttt{cmb\_s4\_errata\_v392}). Lacuna: o parágrafo é de época e trata o CMB-S4 como instrumento esperado, sem datá-lo.

\end{document}
